\documentclass{article}

\usepackage{amsmath,amssymb}
\usepackage{xurl}
\usepackage[hidelinks]{hyperref}
\setlength{\emergencystretch}{2em}

% Shared commands for notes in docs/paper-gaps/.

\DeclareUnicodeCharacter{2080}{\ifmmode{}_{0}\else\textsubscript{0}\fi}
\DeclareUnicodeCharacter{2081}{\ifmmode{}_{1}\else\textsubscript{1}\fi}
\DeclareUnicodeCharacter{2082}{\ifmmode{}_{2}\else\textsubscript{2}\fi}
\DeclareUnicodeCharacter{2099}{\ifmmode{}_{n}\else\textsubscript{n}\fi}

\newcommand{\Lean}{\textsc{Lean}}
\ExplSyntaxOn
\NewDocumentCommand{\leanid}{m}{
  \ifmmode
    \textsf{\detokenize{#1}}
  \else
    \group_begin:
    \urlstyle{sf}
    \tl_set:Nn \l_tmpa_tl {#1}
    \tl_replace_all:Nnn \l_tmpa_tl {\_} {_}
    \exp_args:NV \path \l_tmpa_tl
    \group_end:
  \fi
}
\ExplSyntaxOff
\providecommand{\tr}{\operatorname{tr}}
\newcommand{\tnleanIssueBase}{https://github.com/LionSR/TNLean/issues/}
\newcommand{\tnleanPrBase}{https://github.com/LionSR/TNLean/pull/}
\newcommand{\tnleanDocBase}{https://sirui-lu.com/TNLean/docs/}
\newcommand{\ghissue}[1]{\href{\tnleanIssueBase#1}{GitHub issue \##1}}
\newcommand{\ghpr}[1]{\href{\tnleanPrBase#1}{PR \##1}}
\newcommand{\doclink}[1]{\href{\tnleanDocBase#1.html}{\path{#1}}}

% Dirac notation and the identity operator, as in blueprint/src/macros/common.tex.
% \providecommand keeps note-local definitions authoritative.
\usepackage{dsfont}
\providecommand{\ket}[1]{|#1\rangle}
\providecommand{\bra}[1]{\langle #1|}
\providecommand{\braket}[2]{\langle #1 | #2 \rangle}
\providecommand{\ketbra}[2]{|#1\rangle\!\langle #2|}
\providecommand{\Id}{\mathds{1}}
\providecommand{\one}{\mathds{1}}
\providecommand{\C}{\mathbb{C}}
\providecommand{\MN}[1]{M_{#1}(\C)}
\providecommand{\MD}{M_D(\C)}

% External referencing for paper sources.  The build helper writes sanitized
% auxiliary files under build/paper-aux/ and excludes duplicated source labels.
\usepackage{xr}

% Import a generated paper auxiliary file with a caller-supplied label prefix.
% The candidate paths support builds from docs/paper-gaps/, the repository root,
% and build/paper-aux/ itself.
\newcommand{\importpaperaux}[2]{%
  \IfFileExists{../../build/paper-aux/#2.aux}{%
    \externaldocument[#1]{../../build/paper-aux/#2}%
  }{%
    \IfFileExists{build/paper-aux/#2.aux}{%
      \externaldocument[#1]{build/paper-aux/#2}%
    }{%
      \IfFileExists{#2.aux}{%
        \externaldocument[#1]{#2}%
      }{}%
    }%
  }%
}

% General external reference macros
\newcommand{\paperref}[2]{\ref{#1-#2}}
\newcommand{\papereqref}[2]{(\ref{#1-#2})}

% CPSV16 source (1606.00608 / MPDO-22-12-17-2.tex) external document and shortcuts
\importpaperaux{cpsv16-}{1606.00608/MPDO-22-12-17-2-tnlean}
\newcommand{\cpsvref}[1]{\paperref{cpsv16}{#1}}
\newcommand{\cpsveqref}[1]{\papereqref{cpsv16}{#1}}

% Verdict marker: \gapnote{<kind>}{<status>}. Kind is the policy
% classification (clarification, local-correction, scope-restriction,
% unfaithful, false-source, open-gap); status is open, wip (elimination
% underway), resolved, or historical. Machine-read by the site index;
% prints a verdict line.
\newcommand{\gapnote}[2]{\par\noindent\textbf{Verdict:} #1 (#2).\par}


\title{Trace Preservation in the Forward Direction of Theorem 4.1
  (de las Cuevas--Cirac--Schuch--P\'erez-Garc\'ia)}
\author{}
\date{2026-09-04}

\begin{document}
\maketitle
\gapnote{false-source}{open}

\section{The assertion}

This note concerns the forward implication in Theorem~4.1 of de las Cuevas,
Cirac, Schuch, and P\'erez-Garc\'ia
\cite[Theorem~4.1]{DeLasCuevas2017Irreducible}. The relevant source passages
are the definition of the irreducible forms at
\path{Papers/1708.00029/main.tex:245--266} and
\path{Papers/1708.00029/main.tex:313--331}, the definitions of divisibility
and refinement at \path{Papers/1708.00029/main.tex:717--724}, the theorem at
\path{Papers/1708.00029/main.tex:728--731}, and its forward proof at
\path{Papers/1708.00029/main.tex:735--810}. The companion note
\path{dccsp17_root_kraus_rank_thm41.tex} records the separate gap in the
converse implication.

A tensor $B=(B^1,\ldots,B^d)$ of bond dimension $D$ is in
\emph{irreducible form} when
\begin{align}
    B^i
     & =\bigoplus_{j\in J}\mu_j\,B_j^i,
    \qquad \mu_j>0,
    \label{eq:dccsp17_thm41_forward_trace_preservation_irreducible_form}
\end{align}
with every block map $\mathcal E_{B_j}(X)=\sum_i B_j^iX(B_j^i)^\dagger$
irreducible of spectral radius one. It is in \emph{irreducible form~II} when
moreover every $\mathcal E_{B_j}$ is trace preserving with a diagonal
positive fixed point. The multiplicities $\mu_j$ are arbitrary positive
scalars in both forms; the source normalizes the blocks, not the weights.

A trace-preserving completely positive map $\mathcal E$ is called
\emph{$p$-divisible} when there is a trace-preserving completely positive map
$\mathcal E'$ with $\mathcal E=(\mathcal E')^p$. The family $\mathcal V(B)$
is \emph{$p$-refinable} when there are a tensor $A=(A^1,\ldots,A^d)$ and an
isometry $W:\mathbb C^d\to(\mathbb C^d)^{\otimes p}$ with
\begin{align}
    \ket{V_{pN}(A)}
     & =W^{\otimes N}\ket{V_N(B)}
    \qquad\text{for every }N.
    \label{eq:dccsp17_thm41_forward_trace_preservation_refinement}
\end{align}
Theorem~4.1 asserts, for $B$ in irreducible form~II, that $\mathcal V(B)$ is
$p$-refinable if and only if the transfer map
$\mathcal E_B(X)=\sum_i B^iX(B^i)^\dagger$ is $p$-divisible.

\section{The counterexample}

Let $B_0$ be any tensor in irreducible form~II whose multiplicities all equal
one and whose family is $p$-refinable with witness $(A_0,W)$. The simplest
instance has $d=D=1$, $B_0=(1)$, $A_0=(1)$, and $W=1$, for which
$\ket{V_{pN}(A_0)}=1=\ket{V_N(B_0)}$ at every length. Fix $\lambda>1$ and
set
\begin{align}
    B
     & :=\lambda B_0,
    \qquad
    A
    :=\lambda^{1/p}A_0.
    \label{eq:dccsp17_thm41_forward_trace_preservation_scaled_tensor}
\end{align}
The tensor $B$ is in irreducible form~II: it has the blocks of $B_0$ and the
multiplicities $\lambda\mu_j$, which are positive. The family
$\mathcal V(B)$ is $p$-refinable with witness $(A,W)$, because
\begin{align}
    \ket{V_{pN}(A)}
     & =\lambda^{N}\ket{V_{pN}(A_0)}
    =\lambda^{N}W^{\otimes N}\ket{V_N(B_0)}
    =W^{\otimes N}\ket{V_N(B)}.
    \label{eq:dccsp17_thm41_forward_trace_preservation_scaled_refinement}
\end{align}
Its transfer map is $\mathcal E_B=\lambda^2\mathcal E_{B_0}$, so
$\operatorname{tr}\mathcal E_B(X)=\lambda^2\operatorname{tr}X$. A $p$-th
power of a trace-preserving map is trace preserving, so $\mathcal E_B$ is not
$p$-divisible. The forward implication of Theorem~4.1 therefore fails on
nondegenerate data satisfying its printed hypothesis.

The printed argument breaks at its last line,
\path{Papers/1708.00029/main.tex:806--810}. It constructs $\widetilde A$ with
$\widetilde A^{(p)}=C$ and concludes
$\mathcal E_{\widetilde A}^p=\mathcal E_C=\mathcal E_B$, which is correct,
but $\widetilde A$ is a unitary conjugate of
$A'=\bigoplus_j\mu_j A'_j$, whose transfer map is trace preserving only when
every multiplicity of $A$ equals one. The argument never constrains the
multiplicities of $B$, so it cannot deliver the trace preservation that the
conclusion requires.

\section{The corrected statement}

Add to the hypothesis of the forward implication that $\mathcal E_B$ is trace
preserving. Since
\begin{align}
    \sum_i (B^i)^\dagger B^i
     & =\bigoplus_{j\in J}\mu_j^2\,\mathbf 1_{D_j}
    \label{eq:dccsp17_thm41_forward_trace_preservation_trace_condition}
\end{align}
for a tensor in irreducible form~II, this is the same as requiring every
multiplicity of $B$ to equal one. The added hypothesis is presupposed by the
source: its definition of $p$-divisibility applies only to trace-preserving
maps, so the conclusion ``$\mathcal E_B$ is $p$-divisible'' already asserts
trace preservation of $\mathcal E_B$.

With this hypothesis the printed argument closes. The tensor
$C^{i_1\ldots i_p}=\sum_i W^{(i_1\ldots i_p),i}B^i$ has the blocks and
multiplicities of $B$, all equal to one. The equal-case theorem
\cite[Theorem~3.8]{DeLasCuevas2017Irreducible} matches the blocks of
$A^{(p)}$ with those of $C$ and relates their multiplicities by roots of
unity after reordering. The blocks of $A^{(p)}$ arising from a block $A_j$ of
$A$ carry the multiplicity $\mu_j^p$, so $\mu_j^p$ is a root of unity, and
$\mu_j>0$ forces $\mu_j=1$. Hence $A'$ and $\widetilde A=U^\dagger A'U$ are
left canonical, $\mathcal E_{\widetilde A}$ is trace preserving, and
$\mathcal E_B=\mathcal E_{\widetilde A}^p$ exhibits $p$-divisibility. The
converse implication is unaffected, since it assumes $p$-divisibility and
therefore trace preservation of $\mathcal E_B$ from the start.

\section{The formal statement}

The formal irreducible form is a statement about generated families rather
than about the matrices of the tensor: a tensor is in irreducible form when it
generates, at every length including zero, the same matrix product vectors as
a scalar-weighted direct sum of periodic left-canonical blocks with positive
weights.\footnote{The definition is \leanid{MPSTensor.IsIrreducibleForm} in
    \path{TNLean/MPS/Periodic/Defs.lean}.} It does not fix the transfer map of the
tensor. Two consequences follow for the forward canonicalization hypothesis,
which asks, for every tensor $B$ in this formal irreducible form whose family
is $p$-refinable, for a left-canonical tensor $\widetilde A$ with
$\mathcal E_B=\mathcal E_{\widetilde A^{[p]}}$.\footnote{The hypothesis is
    \leanid{MPSTensor.PRefinementCanonicalization} and the conditional forward
    theorem is \leanid{MPSTensor.thm_4_1_p_refinement_forward}, both in
    \path{TNLean/MPS/Periodic/Symmetry/Theorem41Forward.lean}.}

First, the scaling counterexample refutes the hypothesis at bond dimension
one for every physical dimension and every $p\geq1$: the tensor
$\lambda B_0$ generates the same family as the weighted block sum with weight
$\lambda$, so it is in the formal irreducible form, and the conclusion would
identify the map $x\mapsto\lambda^2x$ with a $p$-th power of the identity.

Second, even after adding trace preservation of $\mathcal E_B$, the formal
irreducible form does not supply the block structure that the printed
argument uses. For a left-canonical normal tensor $B_0$ and an invertible
$Y$, the tensor $YB_0Y^{-1}$ generates the same family as $B_0$ and is in the
formal irreducible form, while
$\mathcal E_{YB_0Y^{-1}}$ is trace preserving exactly when $Y^\dagger Y$ is a
fixed point of $\mathcal E_{B_0}^*$, that is, when $Y$ is a scalar multiple of
a unitary. Recovering a unitary similarity to a block form from equality of
families and trace preservation is a statement about arbitrary tensors that
the cited argument does not contain; the source instead takes $B$ literally
in the form
(\ref{eq:dccsp17_thm41_forward_trace_preservation_irreducible_form}).

The conditional formal theorems remain correct, but their hypothesis cannot
be discharged. The elimination plan is to restate the forward direction for a
tensor given as a multiplicity-bearing block form with every multiplicity
equal to one, the structure already used by the formal equal-case
theorem,\footnote{The equal-case theorem is
    \leanid{MPSTensor.fundamentalTheorem_periodic_equalCase_irreducibleForm} in
    \path{TNLean/MPS/Periodic/EqualCaseGlobal.lean}.} and then to follow the
printed argument. Its ingredients without a formal counterpart are the
blocking lemma for a periodic block,
\path{Papers/1708.00029/main.tex:432--446}, the cyclic projector structure of
a periodic block, \path{Papers/1708.00029/main.tex:335--341}, the unique
decomposition of residues, \path{Papers/1708.00029/main.tex:446--450}, and
the redistribution of the root-of-unity factors across the sites,
\path{Papers/1708.00029/main.tex:765--806}.

\section{Conclusion}

The forward implication of Theorem~4.1 is false as printed: a positive
rescaling of any $p$-refinable tensor in irreducible form~II is again in
irreducible form~II and $p$-refinable, while its transfer map is not trace
preserving and hence not $p$-divisible. The corrected statement adds trace
preservation of $\mathcal E_B$, equivalently unit multiplicities, which the
source's definition of $p$-divisibility presupposes; the printed argument then
proves it. The formal forward canonicalization hypothesis inherits the
counterexample and is refuted independently by non-unitary similarities, so
it must be replaced by a restatement over block-form tensors with
trace-preserving transfer map rather than discharged.

\bibliographystyle{alpha}
\bibliography{references}

\end{document}
